\documentclass[11pt,a4paper]{article}
\usepackage[margin=25mm]{geometry}
\usepackage[T1]{fontenc}
\usepackage[utf8]{inputenc}
\usepackage{lmodern,amsmath,amssymb,enumitem,xurl,xcolor}
\usepackage[colorlinks=true,linkcolor=black,urlcolor=blue!55!black]{hyperref}
\hypersetup{pdfauthor={Luca Scharr},pdftitle={From Images to 3D}}
\setlength{\parindent}{0pt}
\setlength{\parskip}{0.45em}
\setlist{leftmargin=*,itemsep=0.15em,topsep=0.15em}
\emergencystretch=2em
\begin{document}
\begin{center}
{\Large\bfseries Tutorial 10 - From Images to 3D}\\[2mm]
Luca Scharr\\
Computer Vision 2026\\
University of Bern\\[1mm]
Interactive lab: \href{https://sosekie.github.io/Unibe-CVG-CV-Interactive-Labs/tutorial-10/}{Open Tutorial 10 online}
\end{center}
\section*{Learning objectives}
\begin{itemize}[nosep]
\item Derive projection ambiguity and reason about two-view geometry.
\item Explain fitting, learned priors and generalisation from evidence.
\item Compare 3D representations and choose evaluation protocols for a task.
\end{itemize}
Explain your reasoning and state your assumptions.
Use the interactive experiments to check your predictions.
\section*{Part 1. One image and more evidence}
\subsection*{1. One image, many scenes}
\noindent\textbf{Objective.} Derive projection ambiguity and distinguish the image from the physical scene.
\begin{minipage}{\linewidth}
\subsubsection*{1.1}
A point $(X,Y,Z)$ projects to $x=fX/Z$ with a calibrated pinhole camera.
\begin{enumerate}[label=(\alph*),nosep]
\item Multiply all three point coordinates by $k>0$.
What happens to x?
\item Does an unchanged image establish that the physical scene is unchanged?
\end{enumerate}
\end{minipage}
\par\medskip
\begin{minipage}{\linewidth}
\subsubsection*{1.2}
An object is twice as large and twice as far from the camera.
The camera does not move.
\begin{enumerate}[label=(\alph*),nosep]
\item What happens to its projected size under perspective projection?
\item Would a learned estimate of its size make the observation geometrically unique?
\end{enumerate}
\end{minipage}
\par\medskip
\noindent\textbf{Interactive experiment.}
\begin{enumerate}[nosep]
\item Change the object size and distance separately, then lock their scale together.
\item Compare the side view, the image-plane segment and the projected image coordinate.
\item Double the focal length and explain why this is a different change.
\end{enumerate}
\noindent\href{https://sosekie.github.io/Unibe-CVG-CV-Interactive-Labs/tutorial-10/#ambiguity}{Open the interactive experiment.}
\subsection*{2. Combine a prior and a second view}
\noindent\textbf{Objective.} Compare size assumptions and calibrated stereo measurements in the same camera coordinate system.
\begin{minipage}{\linewidth}
\subsubsection*{2.1}
A calibrated camera has focal length f = 1.
An object projects to image height h = 0.5, so $h=fH/Z$, where H is object height and Z is distance.
\begin{enumerate}[label=(\alph*),nosep]
\item Calculate Z for H = 1.5 m and H = 3 m.
Can this image alone distinguish them?
\item Why could a model trained mostly on 1.5 m objects give a plausible but wrong metric estimate for a 3 m object?
\end{enumerate}
\end{minipage}
\par\medskip
\begin{minipage}{\linewidth}
\subsubsection*{2.2}
A toy model assigns probabilities to candidate object heights and uses the fixed image measurement h = 0.5 to obtain candidate distances.
It then receives a disparity measurement from a second calibrated camera.
\begin{enumerate}[label=(\alph*),nosep]
\item What determines the distance distribution before the second measurement?
\item How should the second measurement change the candidate probabilities, and what if that measurement is noisy?
\end{enumerate}
\end{minipage}
\par\medskip
\begin{minipage}{\linewidth}
\subsubsection*{2.3}
Two parallel calibrated cameras with focal length f and baseline B observe a point with disparity $d=u_1-u_2$.
\begin{enumerate}[label=(\alph*),nosep]
\item Derive its depth Z in terms of $f,B,\text{ and }d$.
\item Why does a small baseline make depth sensitive to matching error?
\end{enumerate}
\end{minipage}
\par\medskip
\begin{minipage}{\linewidth}
\subsubsection*{2.4}
A second camera rotates around the same optical centre as the first.
A point is matched accurately across the images.
\begin{enumerate}[label=(\alph*),nosep]
\item Does this rotation alone triangulate the distance of the point?
\item What scale remains ambiguous if camera translation is estimated only from images?
\end{enumerate}
\end{minipage}
\par\medskip
\noindent\textbf{Interactive experiment.}
\begin{enumerate}[nosep]
\item Turn the size prior on and off.
Explain which 3D points can produce the first image.
\item Turn on the second camera and vary the baseline or disparity error.
Compare its ray intersection with the prior-based estimate.
\item Use an unusual object height and discuss when the prior and the second-view measurement disagree.
\end{enumerate}
\begin{minipage}{\linewidth}
\subsection*{Optional extension: Shortcut learning}
A toy model predicts depth as $4+w_{\rm size}(1/h-2)+w_{\rm colour}c$, where h is image height and c = -1 for red or +1 for blue.
The training set contains a red object at 2 m with h = 1 and a blue object at 6 m with h = 1/3.
\begin{enumerate}[label=(\alph*),nosep]
\item Ignoring regularisation, what constraint must the weights satisfy to fit both training examples exactly?
Give two different solutions.
\item Now add a blue object at 2 m and a red object at 6 m.
Which weights fit all four examples exactly, and why does that change generalisation to swapped colours?
\end{enumerate}
\end{minipage}
\par\medskip
\noindent\href{https://sosekie.github.io/Unibe-CVG-CV-Interactive-Labs/tutorial-10/#learning}{Open the interactive experiment.}
\section*{Part 2. Fit a scene}
\subsection*{3. Fit structure and motion}
\noindent\textbf{Objective.} Compute reprojection residuals and explain how joint fitting and incorrect correspondences affect reconstruction.
\begin{minipage}{\linewidth}
\subsubsection*{3.1}
In the experiment, cameras 1 and 2 are at X = 0 and X = 1 with focal length f = 1.
Landmark A has synthetic reference coordinates $(X,Y,Z)=(-0.65,0.42,3.2)$.
The initial estimated coordinates are (-0.42, 0.53, 2.4).
Each camera projects a point to $(u,v)=((X-C)/Z,Y/Z)$, where C is the camera position.
\begin{enumerate}[label=(\alph*),nosep]
\item Compute the observed and initial projected coordinates for A in each of the first two images, and compute the mean of the two squared 2D residuals.
\item The fitting control also adjusts the other three points and camera 3.
Why must these parameters be considered together in the full objective?
\end{enumerate}
\end{minipage}
\par\medskip
\begin{minipage}{\linewidth}
\subsubsection*{3.2}
All four landmarks are visible in all three images, giving 12 two-dimensional observations.
The experiment corrupts landmark C in camera 3 by adding 0.22 to its horizontal observation and subtracting 0.12 from its vertical observation.
\begin{enumerate}[label=(\alph*),nosep]
\item If the estimated geometry initially equals the true geometry, calculate the corrupted observation's squared 2D residual and the mean squared residual across all 12 observations.
\item What can happen when the fitting algorithm tries to reduce this error, and how could one limit the influence of the wrong match?
\end{enumerate}
\end{minipage}
\par\medskip
\noindent\textbf{Interactive experiment.}
\begin{enumerate}[nosep]
\item Select landmark A.
Compare its filled observations and outlined projections in cameras 1 and 2; record the mean reprojection error.
\item Take one fitting step, then refine the reconstruction.
Observe how all four point estimates and the third camera can change together.
\item Reset, corrupt landmark C in camera 3, and refine again.
Reveal the synthetic reference and compare the geometry with the image residuals.
\end{enumerate}
\noindent\textbf{Optional SfM factorisation.} Why does centering remove the image translations?
Why does any invertible $3\times 3$ matrix leave the factorisation ambiguous?
\noindent\href{https://sosekie.github.io/Unibe-CVG-CV-Interactive-Labs/tutorial-10/#sfm}{Open the interactive experiment.}
\section*{Part 3. Representations and evaluation}
\subsection*{4. Choose a 3D representation}
\noindent\textbf{Objective.} Compare what depth, point maps, occupancy and Gaussian outputs encode.
\begin{minipage}{\linewidth}
\subsubsection*{4.1}
Choose a representation for measuring visible distances, building a closed object mesh, and rendering a nearby novel view.
\begin{enumerate}[label=(\alph*),nosep]
\item Which representation is suitable for each task?
\item What limitation or validation requirement accompanies each choice?
\end{enumerate}
\end{minipage}
\par\medskip
\begin{minipage}{\linewidth}
\subsubsection*{4.2}
A network predicts a depth z and a camera ray direction r for one pixel.
\begin{enumerate}[label=(\alph*),nosep]
\item How do the direction convention and the definition of depth determine the 3D point?
\item Why can two plausible predicted hidden surfaces be consistent with the same visible image?
\end{enumerate}
\end{minipage}
\par\medskip
\noindent\textbf{Interactive experiment.}
\begin{enumerate}[nosep]
\item Choose a representation and inspect the reference pole, its hidden back, the empty space and the wall.
\item Compare the visible-surface depth and point maps with the occupancy field and Gaussian primitives.
\item Explain which information can be represented and what would require additional evidence or a prior.
\end{enumerate}
\noindent\href{https://sosekie.github.io/Unibe-CVG-CV-Interactive-Labs/tutorial-10/#representations}{Open the interactive experiment.}
\clearpage
\subsection*{5. Evaluate depth and geometry}
\noindent\textbf{Objective.} Explain the effect of scale alignment and why averages can hide local depth failures.
\begin{minipage}{\linewidth}
\subsubsection*{5.1}
Prediction A is exactly twice the reference depth at every pixel.
Prediction B predicts the wall depth for every pixel and therefore misses the narrow foreground pole.
\begin{enumerate}[label=(\alph*),nosep]
\item What happens to A's error if a best-fit global scale is permitted?
\item Why can B have a small whole-image error while missing an important object?
\end{enumerate}
\end{minipage}
\par\medskip
\begin{minipage}{\linewidth}
\subsubsection*{5.2}
A Gaussian reconstruction matches its training photographs well and yields a plausible held-out rendering.
\begin{enumerate}[label=(\alph*),nosep]
\item What evidence does the held-out rendering provide?
\item Does that alone validate a collision surface in hidden regions?
\end{enumerate}
\end{minipage}
\par\medskip
\noindent\textbf{Interactive experiment.}
\begin{enumerate}[nosep]
\item In comparison 1, inspect the original prediction and fit a global scale.
Compare the before and after depth errors and explain what each score measures.
\item In comparison 2, vary the predicted pole depth.
Compare the pole-region MSE with the whole-image MSE when the pole disappears.
\item Explain why a held-out view tests a different claim from an independent geometry measurement.
\end{enumerate}
\noindent\href{https://sosekie.github.io/Unibe-CVG-CV-Interactive-Labs/tutorial-10/#evaluation}{Open the interactive experiment.}
\end{document}
